\section{Introdução}

Este documento apresenta uma sugestão de modelo para os Trabalhos de Conclusão de
Curso de Engenharia de Controle e Automação de Processos. Sendo uma sugestão, as
seções podem ser reestruturadas (inclusão, exclusão, definição de nomes das seções e
subseções) para melhor adequação do trabalho, principalmente, em função dos diferentes
formatos de apresentação dos Trabalhos de Conclusão de Curso. Por exemplo, para
trabalhos apresentados no formato de software, uma possível estrutura é: (i) Introdução,
(ii) Fundamentação Teórica, (iii) Apresentação do software, (iv) Guia de Uso, (v)
Referências, (vi) Apêndices.
Esta primeira seção é voltada à descrever a contextualização e motivação do seu
trabalho, incorporando a justificativa, e associando-as ao objetivo geral e objetivos
específicos.

Utilize referências atualizadas, preferencialmente dos últimos 5 anos, para um
embasamento do seu trabalho.
Todo conteúdo literal ou ideia extraído de um documento publicado, deve ser
devidamente citado. Na citação indireta a referência ao documento ou autor de deve ser
feita entre parênteses e indicar o ano da publicação, como por exemplo (Guerrieri, 2017).
Esse tipo de citação é feito ao longo ou o final de um parágrafo.
Quando a citação é feita de forma direta, apenas o ano da publicação deve vir entre
parênteses. Por exemplo, “conforme Krenz (2005) ...”. Caso o trabalho tenha sido escrito
por 3 ou mais autores, apenas o primeiro deve ser citado e o termo et al deve ser
adicionado (Donati et al, 1981; Costa et al, 2010). O MS Word possui um recurso para
gerenciar isso, caso desejem usar.
Recomendo utilizar o sistema autor-data da ABNT para citação, consulte Lubisco
e Viera (2019).
1.1
Objetivo Geral
Se desejar, pode incluir uma subseção. Formate-a com o “estilo” “Título 2”.